% \iffalse meta-comment
% A LaTeX Class for Jacobs University Slides
% Copyright (c) 2007 Michael Kohlhase, all rights reserved
%               this file is released under the
%               Gnu Library Public Licences (LGPL)
%
% The original of this file is in the public repository at 
% http://github.com/KWARC/sTeX/
% \fi
% 
% \iffalse
%<cls|package>\NeedsTeXFormat{LaTeX2e}[1999/12/01]
%<cls>\ProvidesClass{mikoslides}[2016/04/07 v1.1 MiKo slides Class]
%<package>\ProvidesPackage{mikoslides}[2016/04/07 v1.1 MiKo slides Package]
%
%<*driver>
\documentclass{ltxdoc}
\usepackage[utf8]{inputenc}
\usepackage{url,array,float,xspace}
\usepackage[show]{ed}
\usepackage{graphicx,paralist}
\usepackage[hyperref=auto,style=alphabetic]{biblatex}
\addbibresource{kwarcpubs.bib}
\addbibresource{extpubs.bib}
\addbibresource{kwarccrossrefs.bib}
\addbibresource{extcrossrefs.bib}
\makeindex
\floatstyle{boxed}
\newfloat{exfig}{thp}{lop}
\floatname{exfig}{Example}
\def\githubissue#1{\cite{sTeX:github:on},
\hyperlink{https://github.com/KWARC/sTeX/issues/#1}{issue #1}}
\usepackage{hyperref}
\begin{document}
\RecordChanges
\DocInput{mikoslides.dtx}
\end{document}
%</driver>
% \fi
% 
% \iffalse\CheckSum{469}\fi
% 
% \changes{v0.1}{2005/12/06}{Initial Version}
% \changes{v0.2}{2006/01/11}{course notes back on seminar}
% \changes{v0.3}{2007/05/16}{changing to Jacobs logo}
% \changes{v0.4}{2007/10/16}{re-basing the whole thing on beamer}
% \changes{v0.4}{2008/09/06}{moving line-end-comment to \texttt{omdoc.dtx}}
% \changes{v0.5}{2009/06/17}{eliminating mytwocolumns, this is better done by \texttt{beamer.cls}}
% \changes{v0.9}{2010/06/18}{this is almost done}
% \changes{v0.9}{2012/09/17}{basic options handling for the \texttt{frame} environment in
% notes mode}
% \changes{v0.9}{2013/08/23}{numbered sectocframes}
% \changes{v1.0}{2014/01/07}{adding \texttt{\textbackslash frameimage}}
% \changes{v1.1}{2015/10/25}{Removing the old title macros (use the regular ones instead)}
% \changes{v1.1}{2015/10/25}{reinterpreting omgroup}
% \changes{v1.1}{2015/11/04}{moving MathHub support out to separate package}
% 
% \GetFileInfo{mikoslides.cls}
% \MakeShortVerb{\|}
%
% \def\twin#1#2{\index{#1!#2}\index{#2!#1}}
% \def\twintoo#1#2{{#1 #2}\twin{#1}{#2}}
% \def\atwin#1#2#3{\index{#1!#2!#3}\index{#3!#2 (#1)}}
% \def\atwintoo#1#2#3{{#1 #2 #3}\atwin{#1}{#2}{#3}}
%
% \def\scsys#1{{{\sc #1}}\index{#1@{\sc #1}}}
% \def\stex{\hbox{\raisebox{-.5ex}S\kern-.5ex\TeX}}
% \def\sTeX{\stex}
% \def\cnxml{\scshape{CNXml}}
% \def\connexions{\scshape{Connexions}}
% \def\element#1{{\ttfamily{#1}}}
% \def\snippet#1{{\ttfamily{#1}}}
% \def\cnxlatex{CNX\LaTeX\xspace}
% \def\mathml{{\scshape{MathML}}\xspace}
% \def\omdoc{OMDoc\xspace}
% \def\activemath{{\scshape{ActiveMath}}\xspace}
% \def\textwarning{\includegraphics[width=1.2em]{dangerous-bend}\xspace}
% 
% \title{Slides and Course Notes\thanks{Version {\fileversion}
% (last revised {\filedate})}}
%    \author{Michael Kohlhase\\
%            Jacobs University, Bremen\\
%            \url{http://kwarc.info/kohlhase}}
% \maketitle
%
% \begin{abstract}
%   We present a document class from which we can generate both course slides and course
%   notes in a transparent way.
% \end{abstract}
%
% \tableofcontents\newpage
%
%\section{Introduction}
%
% The |mikoslides| document class is derived from |beamer.cls|~\cite{beamerclass:on}, it
% adds a ``notes version'' for course notes derived from the |omdoc|
% class~\cite{Kohlhase:smomdl} that is more suited to printing than the one supplied by
% |beamer.cls|.
%
% 
%\section{The User Interface}\label{sec:user}
%
% The |mikoslides| class takes the notion of a slide frame from Till Tantau's excellent
% |beamer| class and adapts its notion of frames for use in the \sTeX and \omdoc. To
% support semantic course notes, it extends the notion of mixing frames and explanatory
% text, but rather than treating the frames as images (or integrating their contents into
% the flowing text), the |mikoslides| package displays the slides as such in the course
% notes to give students a visual anchor into the slide presentation in the course (and to
% distinguish the different writing styles in slides and course notes).
% 
% In practice we want to generate two documents from the same source: the slides for
% presentation in the lecture and the course notes as a narrative document for home
% study. To achieve this, the |mikoslides| class has two modes: \emph{slides mode} and
% \emph{notes mode} which are determined by the package option. 
%
% \subsection{Package Options}\label{sec:user:options}
% 
% The |mikoslides| class takes a variety of class options:\ednote{leaving out noproblems
% for the moment until we decide what to do with it.}
% \begin{itemize}
% \item The options \DescribeMacro{slides}|slides|\DescribeMacro and {notes}|notes| switch
%   between slides mode and notes mode (see Section~\ref{sec:user:notesslides}).
% \item If the option \DescribeMacro{sectocframes}|sectocframes| is given, then special
%   frames with section table of contents are produced headers \ednote{document the functionality}
% \item \DescribeMacro{showmeta}|showmeta|. If this is set, then the metadata keys are
%   shown (see~\cite{Kohlhase:metakeys:ctan} for details and customization options).
% \item If the option \DescribeMacro{frameimages}|frameimages| is set, then slide mode
%   also shows the |\frameimage|-generated frames. 
% \end{itemize}
% 
% \subsection{Notes and Slides}\label{sec:user:notesslides}
% 
% Slides are represented with the \DescribeEnv{frame}|frame| just like in the |beamer|
% class, see~\cite{Tantau:ugbc} for details. The |mikoslides| class adds the
% \DescribeEnv{note}|note| environment for encapsulating the course note
% fragments.\footnote{MK: it would be very nice, if we did not need this environment, and
% this should be possible in principle, but not without intensive LaTeX trickery. Hints to
% the author are welcome.} 
%
% \textwarning Note that it is essential to start and end the |notes| environment at the
% start of the line -- in particular, there may not be leading blanks -- else {\LaTeX}
% becomes confused and throws error messages that are difficult to decipher.
%
% \begin{exfig}
% \begin{verbatim}
% \begin{note}
%   We start this course with ...
% \end{note}
%
% \begin{frame}
%   \frametitle{The first slide}
%   ...
% \end{frame}
% \begin{note}
%   ... and more explanatory text
% \end{note}
%
% \begin{frame}
%   \frametitle{The second slide}
%   ...
% \end{frame}
% ...
% \end{verbatim}
% \caption{A typical Course Notes File}\label{fig:notesfile}
% \end{exfig}
% 
% By interleaving the |frame| and |note| environments, we can build course notes as shown
% in Figure~\ref{fig:notesfile}. 
% 
% Sometimes, we want to integrate slides as images after all -- e.g. because we already
% have a PowerPoint presentation, to which we want to add \sTeX notes. In this case we can
% use \DescribeMacro{\frameimage}|\frameimage[|\meta{opt}|]{|\meta{path}|}|, where
% \meta{opt} are the options of |\includegraphics| from the |graphicx|
% package~\cite{CarRah:tpp99} and \meta{path} is the file path (extension can be left off
% like in |\includegraphics|).
% 
% If we want to transclude a the contents of a file as a note, we can use the
% \DescribeMacro{\ninputref}|\ninputref| macro. |\ninputref{foo}| is equivalent to 
%\begin{verbatim}
% \begin{note}
% \inputref{foo}
% \end{note}
% \end{verbatim}
%
% \subsection{Header and Footer Lines}\label{sec:user:headfootlines}
%
% \subsection{Colors and Highlighting}\label{sec:user:highlighting}
% The \DescribeMacro{\textwarning}|\textwarning| macro generates a warning
% sign: \textwarning
%
% \subsection{Front Matter, Titles, etc}\label{sec:user:matter}
%
% \subsection{Miscellaneous}\label{sec:user:misc}
%
% \section{Limitations}\label{sec:limitations}
% 
% In this section we document known limitations. If you want to help alleviate them,
% please feel free to contact the package author. Some of them are currently discussed in
% the \sTeX GitHub repository~\cite{sTeX:github:on}. 
% \begin{enumerate}
% \item when option |book| which uses |\pagestyle{headings}| is given and semantic macros
%   are given in the |omgroup| titles, then they sometimes are not defined by the time the
%   heading is formatted. Need to look into how the headings are made. This is a problem
%   of the underlying |omdoc| package.
% \end{enumerate}
% 
% \StopEventually{\newpage\PrintIndex\newpage\PrintChanges\printbibliography}
% 
%\section{The Implementation}\label{sec:impl}
%
%\subsection{Class and Package Options}\label{sec:impl:init}
%
% We define some Package Options and switches for the |mikoslides| class and activate them
% by passing them on to |beamer.cls| and |omdoc.cls| and the |mikoslides| package. We pass
% the |nontheorem| option to the |statements| package when we are not in notes mode, since
% the |beamer| package has its own (overlay-aware) theorem environments. 
%
%    \begin{macrocode}
%<*cls>
\newif\ifnotes\notesfalse
\DeclareOption{notes}{\notestrue\PassOptionsToPackage{\CurrentOption}{mikoslides}}
\DeclareOption{slides}{\notesfalse\PassOptionsToPackage{\CurrentOption}{mikoslides}}
\DeclareOption*{\PassOptionsToClass{\CurrentOption}{omdoc}
                           \PassOptionsToClass{\CurrentOption}{beamer}
                           \PassOptionsToPackage{\CurrentOption}{mikoslides}}
\ProcessOptions
%</cls>
%    \end{macrocode}
% now we do the same for the |mikoslides| package. Note that we also have to define the
% same switches\ednote{MK: we may think about making all of them internal}, since we might
% use |mikoslides.sty| in a different class.
%    \begin{macrocode}
%<*package>
\newif\if@mikoslides@mh@\@mikoslides@mh@false
\DeclareOption{mh}{\@mikoslides@mh@true
  \PassOptionsToPackage{\CurrentOption}{stex}
  \PassOptionsToPackage{\CurrentOption}{smglom}
  \PassOptionsToPackage{\CurrentOption}{tikzinput}}
\newif\ifnotes\notesfalse
\DeclareOption{notes}{\notestrue}
\DeclareOption{slides}{\notesfalse}
\ifnotes\else\PassOptionsToPackage{nontheorem}{statements}\fi
\newif\ifsectocframes\sectocframesfalse
\DeclareOption{sectocframes}{\sectocframestrue\PassOptionsToPackage{msection}{statements}}
\newif\ifframeimages\frameimagesfalse
\DeclareOption{frameimages}{\frameimagestrue}
\newif\if@part\@partfalse
\DeclareOption{report}{\@parttrue\PassOptionsToPackage{\CurrentOption}{omdoc}}
\DeclareOption{book}{\@parttrue\PassOptionsToPackage{\CurrentOption}{omdoc}}
\DeclareOption{bookpart}{\@parttrue\PassOptionsToPackage{\CurrentOption}{omdoc}}
\newif\ifproblems\problemstrue
\DeclareOption{noproblems}{\problemsfalse}
\DeclareOption*{\PassOptionsToPackage{\CurrentOption}{stex}
                           \PassOptionsToPackage{\CurrentOption}{smglom}
                           \PassOptionsToPackage{\CurrentOption}{tikzinput}}
\ProcessOptions
%</package>
%    \end{macrocode}
%
% Depending on the options, we either load the |article|-based |omdoc| or the |beamer|
% class (and set some counters).
%    \begin{macrocode}
%<*cls>
\ifnotes
  \LoadClass{omdoc}
\else
  \LoadClass[10pt,notheorems]{beamer}
  \newcounter{Item}
  \newcounter{paragraph}
  \newcounter{subparagraph}
  \newcounter{Hfootnote}
\fi
%    \end{macrocode}
% now it only remains to load the |mikoslides| package that does all the rest. 
%    \begin{macrocode}
\RequirePackage{mikoslides}
%</cls>
%    \end{macrocode}
% 
% In |notes| mode, we also have to make the |beamer|-specific things available to
% |article| via the |beamerarticle| package. We use options to avoid loading theorem-like
% environments, since we want to use our own from the $\sTeX$ packages.  The first batch
% of packages we want are loaded on |mikoslides.sty|. These are the general ones, we will
% load the \sTeX-specific ones after we have done some work (e.g. defined the counters
% |m*|). Only the |stex-logo| package is already needed now for the default theme.
%
%    \begin{macrocode}
%<*package>
\ifnotes
\RequirePackage{a4wide}
\RequirePackage{marginnote}
\RequirePackage{xcolor}
\RequirePackage{mdframed}
\RequirePackage[noxcolor,noamsthm]{beamerarticle}
\fi
\if@mikoslides@mh@\RequirePackage{mikoslides-mh}\fi
\RequirePackage{etoolbox}
\RequirePackage{amssymb}
\RequirePackage{amsmath}
\RequirePackage{comment}
\RequirePackage{textcomp}
\RequirePackage{url}
\RequirePackage{graphicx}
\RequirePackage{stex-logo}
\RequirePackage{pgf}
\ifnotes
\RequirePackage[bookmarks,bookmarksopen,bookmarksnumbered,breaklinks,
linkcolor=black,citecolor=black,urlcolor=cyan,filecolor=cyan,colorlinks]{hyperref}
\fi
%    \end{macrocode}
%
% finally, we require the |metakeys| package from \sTeX, so that we can use the
% |\addmetakey| mechanism.
%
%    \begin{macrocode}
\RequirePackage{metakeys}
%    \end{macrocode}
% 
% \subsection{Notes and Slides}\label{sec:impl:noteslides}
%
% For the lecture notes cases, we also provide the |\usetheme| macro that would otherwise
% from the the |beamer| class. While the latter loads |beamertheme|\meta{theme}{.sty}, the
% notes version loads |beamernotestheme|\meta{theme}|.sty|.\ednote{MK: This is not ideal, but
% I am not sure that I want to be able to provide the full theme functionality there.}
%    \begin{macrocode}
\ifnotes
\renewcommand\usetheme[2][]{\usepackage[#1]{beamernotestheme#2}}
\fi
%    \end{macrocode}
% We define the sizes of slides in the notes. Somehow, we cannot get by with the same
% here. 
%
%    \begin{macrocode}
\newcounter{slide}
\newlength{\slidewidth}\setlength{\slidewidth}{12.8cm}
\newlength{\slideheight}\setlength{\slideheight}{9cm}
%    \end{macrocode}
% 
% \begin{environment}{note}
% The |note| environment is used to leave out text in the |slides| mode. It does not have
% a counterpart in OMDoc. So for course notes, we define the |note| environment to be a
% no-operation otherwise we declare the |note| environment as a comment via the |comment|
% package.
%    \begin{macrocode}
\ifnotes%
  \renewenvironment{note}{\ignorespaces}{}%
\else%
  \excludecomment{note}%
\fi%
%    \end{macrocode}
% \end{environment}
% 
% \begin{macro}{\ninputref}
%    \begin{macrocode}
\newcommand\ninputref[2][]{\ifnotes\inputref[#1]{#2}\fi}
%    \end{macrocode}
% \end{macro}
% 
%
% We first set up the slide boxes in |article| mode. We set up sizes and provide a
% box register for the frames and a counter for the slides.
% 
%    \begin{macrocode}
\ifnotes
  \newlength{\slideframewidth}
  \setlength{\slideframewidth}{1.5pt}
%    \end{macrocode}
% 
% \begin{environment}{frame}
%   We first define the keys. 
%    \begin{macrocode}
  \addmetakey{frame}{label}
  \addmetakey[yes]{frame}{allowframebreaks}
  \addmetakey{frame}{allowdisplaybreaks}
  \addmetakey[yes]{frame}{fragile}
  \addmetakey[yes]{frame}{shrink}
  \addmetakey[yes]{frame}{squeeze}
%    \end{macrocode}
% We define the environment, read them, and construct the slide number and label.
%    \begin{macrocode}
  \renewenvironment{frame}[1][]{%
    \metasetkeys{frame}{#1}%
    \stepcounter{slide}%
    \def\@currentlabel{\theslide}%
    \ifx\frame@label\@empty%
    \else%
      \label{\frame@label}%
    \fi%
%    \end{macrocode}
%   We redefine the |itemize| environment so that it looks more like the one in |beamer|. 
%    \begin{macrocode}
    \def\itemize@level{outer}%
    \def\itemize@outer{outer}%
    \def\itemize@inner{inner}%
    \renewcommand\newpage{}%
    \renewcommand\metakeys@show@keys[2]{\marginnote{{\scriptsize ##2}}}%
    \renewenvironment{itemize}{%
      \ifx\itemize@level\itemize@outer%
        \def\itemize@label{$\rhd$}%
      \fi%
      \ifx\itemize@level\itemize@inner%
        \def\itemize@label{$\scriptstyle\rhd$}%
      \fi%
      \begin{list}%
      {\itemize@label}%
      {\setlength{\labelsep}{.3em}%
       \setlength{\labelwidth}{.5em}%
       \setlength{\leftmargin}{1.5em}%
      }%
      \edef\itemize@level{\itemize@inner}%
    }{%
      \end{list}%
    }%
%    \end{macrocode}
% We create the box with the |mdframed| environment from the equinymous package.
%    \begin{macrocode}
    \begin{mdframed}[linewidth=\slideframewidth,skipabove=1ex,skipbelow=1ex,userdefinedwidth=\slidewidth,align=center]\sf%
  }{%
    \medskip\miko@slidelabel\end{mdframed}%
  }%
%    \end{macrocode}
% \end{environment}
% 
% Now, we need to redefine the frametitle (we are still in course notes mode). 
% \begin{macro}{\frametitle}
%    \begin{macrocode}
  \renewcommand{\frametitle}[1]{{\Large\bf\sf\color{blue}{#1}}\medskip}%
\fi %ifnotes
%    \end{macrocode}
% \end{macro}
% 
% \begin{macro}{\frameimage}
%   We have to make sure that the width is overwritten, for that we check the
%   |\Gin@ewidth| macro from the |graphicx| package\ednote{MK@DG; we need to do that in
%   the LaTeXML binding as well!}
%    \begin{macrocode}
\newrobustcmd\frameimage[2][]{%
  \stepcounter{slide}%
  \ifframeimages%
    \def\Gin@ewidth{}\setkeys{Gin}{#1}%
    \ifnotes%
    \else%
      \vfill%
    \fi%
    \ifx\Gin@ewidth\@empty%
      \mycgraphics[width=\slidewidth,#1]{#2}\else\mycgraphics[#1]{#2}%
    \fi%
    \par\strut\hfill{\footnotesize Slide \arabic{slide}}%
    \ifnotes%
    \else%
      \vfill%
    \fi%
  \fi%
}% ifframeimages
%    \end{macrocode}
% \end{macro}
% 
% \begin{macro}{\pause}
%   \ednote{MK: fake it in notes mode for now}
%    \begin{macrocode}
\ifnotes\newcommand\pause{}\fi
%    \end{macrocode}
% \end{macro}
% 
% \subsection{Header and Footer Lines}\label{sec:impl:headfootlines}
%
% Now, we set up the infrastructure for the footer line of the slides, we use boxes for
% the logos, so that they are only loaded once, that considerably speeds up processing.
% 
% \begin{macro}{\setslidelogo}
% The default logo is the logo of Jacobs University. Customization can be done by |\setslidelogo{|\meta{logo name}|}|.
%    \begin{macrocode}
\newlength{\slidelogoheight}
\ifnotes%
  \setlength{\slidelogoheight}{.4cm}%
\else%
  \setlength{\slidelogoheight}{1cm}%
\fi%
\newsavebox{\slidelogo}%
\sbox{\slidelogo}{\sTeX}%
\newrobustcmd{\setslidelogo}[1]{%
  \sbox{\slidelogo}{\includegraphics[height=\slidelogoheight]{#1}}%
}%
%    \end{macrocode}
% \end{macro}
%
% \begin{macro}{\setsource}
% |\source| stores the writer's name. By default it is {\it Michael Kohlhase} since he is the main user and designer of this package. |\setsource{|\meta{name}|}| can change the writer's name.
%    \begin{macrocode}
\def\source{Michael Kohlhase}% customize locally
\newrobustcmd{\setsource}[1]{\def\source{#1}}%
%    \end{macrocode}
% \end{macro}
%
% \begin{macro}{\setlicensing}
%   Now, we set up the copyright and licensing. By default we use the Creative Commons
%   Attribuition-ShareAlike license to strengthen the public domain. If package |hyperref|
%   is loaded, then we can attach a hyperlink to the license
%   logo. |\setlicensing[|\meta{url}|]{|\meta{logo name}|}| is used for customization,
%   where ||\meta{url}|| is optional.
%    \begin{macrocode}
\def\copyrightnotice{\footnotesize\copyright:\hspace{.3ex}{\source}}%
\newsavebox{\cclogo}%
\sbox{\cclogo}{\includegraphics[height=\slidelogoheight]{cc_somerights}}%
\newif\ifcchref\cchreffalse%
\AtBeginDocument{%
  \@ifpackageloaded{hyperref}{\cchreftrue}{\cchreffalse}
}%
\def\licensing{%
  \ifcchref%
    \href{http://creativecommons.org/licenses/by-sa/2.5/}{\usebox{\cclogo}}%
  \else%
    {\usebox{\cclogo}}%
  \fi%
}%
\newrobustcmd{\setlicensing}[2][]{%
  \def\@url{#1}%
  \sbox{\cclogo}{\includegraphics[height=\slidelogoheight]{#2}}%
  \ifx\@url\@empty%
    \def\licensing{{\usebox{\cclogo}}}%
  \else%
    \def\licensing{%
      \ifcchref%
      \href{#1}{\usebox{\cclogo}}%
      \else%
      {\usebox{\cclogo}}%
      \fi%
    }%
  \fi%
}%
%    \end{macrocode}
% \end{macro} 
%
% \begin{macro}{\slidelabel}
% Now, we set up the slide label for the |article| mode.\ednote{see that we can use the themes for the slides some day. This is all fake.}
%    \begin{macrocode}
\newrobustcmd\miko@slidelabel{%
  \vbox to \slidelogoheight{%
    \vss\hbox to \slidewidth%
    {\licensing\hfill\copyrightnotice\hfill\arabic{slide}\hfill\usebox{\slidelogo}}%
  }%
}%
%    \end{macrocode}
% \end{macro}
% 
% \subsection{Colors and Highlighting}\label{sec:impl:highlighting}
%
% We first specify sans serif fonts as the default. 
%
%    \begin{macrocode}
\sffamily
%    \end{macrocode}
%
% Now, we set up an infrastructure for highlighting phrases in slides. Note that we use
% content-oriented macros for highlighting rather than directly using color markup. 
% The first thing to to is to adapt the green so that it is dark enough for most beamers
%    \begin{macrocode}
\AtBeginDocument{%
\definecolor{green}{rgb}{0,.5,0}%
\definecolor{purple}{cmyk}{.3,1,0,.17}%
}%
%    \end{macrocode}
%
% We customize the |\defemph|, |\notemph|, and |\stDMemph| macros with colors for the use
% in the |statements| package. Furthermore we customize the |\@@lec| macro for the
% appearance of line end comments in |\lec|.
%
%    \begin{macrocode}
% \def\STpresent#1{\textcolor{blue}{#1}}
\def\defemph#1{{\textcolor{magenta}{#1}}}
\def\notemph#1{{\textcolor{magenta}{#1}}}
\def\stDMemph#1{{\textcolor{blue}{#1}}}
\def\@@lec#1{(\textcolor{green}{#1})}
%    \end{macrocode}
%
% I like to use the dangerous bend symbol for warnings, so we provide it here.
% \begin{macro}{\textwarning}
%   as the macro can be used quite often we put it into a box register, so that it is only
%   loaded once. 
%    \begin{macrocode}
\pgfdeclareimage[width=.9em]{miko@small@dbend}{dangerous-bend}
\def\smalltextwarning{%
  \pgfuseimage{miko@small@dbend}%
  \xspace%
}%
\pgfdeclareimage[width=1.5em]{miko@dbend}{dangerous-bend}
\newrobustcmd\textwarning{%
  \raisebox{-.05cm}{\pgfuseimage{miko@dbend}}%
  \xspace%
}%
\pgfdeclareimage[width=2.5em]{miko@big@dbend}{dangerous-bend}%
\newrobustcmd\bigtextwarning{%
  \raisebox{-.05cm}{\pgfuseimage{miko@big@dbend}}%
  \xspace%
}%
%    \end{macrocode}
% \end{macro}
% 
%    \begin{macrocode}
\newrobustcmd\putgraphicsat[3]{%
  \begin{picture}(0,0)\put(#1){\includegraphics[#2]{#3}}\end{picture}%
}%
\newrobustcmd\putat[2]{%
  \begin{picture}(0,0)\put(#1){#2}\end{picture}%
}%
%    \end{macrocode}
%
% \subsection{Sectioning}
%
% If the |sectocframes| option is set, then we make section frames. We first define a set
% of counters\ednote{I forget: why not use the counters from beamer/article? --> document
% this.}
%    \begin{macrocode}
\ifsectocframes%
\if@part%
    \newcounter{mpart}%
    \newcounter{mchapter}%
    \newcounter{msection}[mchapter]%
  \else%
    \newcounter{msection}%
  \fi%
  \newcounter{msubsection}[msection]%
  \newcounter{msubsubsection}[msubsection]%
  \newcounter{msubsubsubsection}[msubsubsection]%
\fi% ifsectocframes
%    \end{macrocode}
%
% Now that we have defined the counters, we can load the \sTeX-specific packages (in
% particular |statements| that needs these counters). 
%
%    \begin{macrocode}
\RequirePackage{stex}
\RequirePackage{smglom}
\RequirePackage{tikzinput}
%    \end{macrocode}
% 
% Now \sTeX is loaded, we redefine the |omgroup| environment to produce section toc frames
% (if the option |sectocframes| is specified.)\ednote{MK: we should probably just redefine
% omgroup@num and omgroup@nonum, since they do the actual work so that we can add the
% sectocframes behavior here without having to copy the internals. Then there is less
% material that can get out of sync. Additionally, we should have a hook in the original
% code of those so that we can increment the slides counter in notes node (to keep slides
% in sync)}
% 
%    \begin{macrocode}
\ifnotes\else% only in slides
  \renewenvironment{omgroup}[2][]{%
    \metasetkeys{omgroup}{#1}\sref@target%
    \advance\section@level by 1%
    \ifsectocframes%
    \stepcounter{slide}
    \begin{frame}[noframenumbering]%
    \vfill\Large\centering%
    \red{%
      \ifcase\section@level\or
      \stepcounter{mpart}
      \def\@@label{Part \Roman{mpart}}
      \def\currentsectionlevel{part}
      \or%
      \stepcounter{mchapter}
      \def\@@label{Chapter \arabic{mchapter}}
      \def\currentsectionlevel{chapter}
      \or
      \stepcounter{msection}
      \def\@@label{\if@part\arabic{mchapter}.\fi\arabic{msection}}
      \def\currentsectionlevel{section}
      \or
      \stepcounter{msubsection}
      \def\@@label{\if@part\arabic{mchapter}.\fi\arabic{msection}.\arabic{msubsection}}
      \def\currentsectionlevel{subsection}
      \or
      \stepcounter{msubsubsection}
      \def\@@label{\if@part\arabic{mchapter}.\fi\arabic{msection}.\arabic{msubsection}.\arabic{msubsubsection}}
      \def\currentsectionlevel{subsubsection}
      \or
      \stepcounter{msubsubsubsection}
      \def\@@label{\if@part\arabic{mchapter}.\fi\arabic{msection}.\arabic{msubsection}.\arabic{msubsubsection}.\arabic{msubsubsubsection}}
      \def\currentsectionlevel{subsubsubsection}
      \fi% end ifcase
      \@@label\sref@label@id\@@label
      \quad #2%
    }%
    \vfill%
    \end{frame}%
    \fi %ifsectocframes
  }
  {\advance\section@level by -1}%
\fi% ifnotes
%    \end{macrocode}
%
% \subsection{Miscellaneous}
%
% We set up a |beamer| template for theorems like ams style, but without a block
% environment.  
%    \begin{macrocode}
\def\inserttheorembodyfont{\normalfont}
\defbeamertemplate{theorem begin}{miko}
{\inserttheoremheadfont\inserttheoremname\inserttheoremnumber
  \ifx\inserttheoremaddition\@empty\else\ (\inserttheoremaddition)\fi%
  \inserttheorempunctuation\inserttheorembodyfont\xspace}
\defbeamertemplate{theorem end}{miko}{}
%    \end{macrocode}
% and we set it as the default one. 
%    \begin{macrocode}
\setbeamertemplate{theorems}[miko]
%    \end{macrocode}
% The following fixes an error I do not understand, this has something to do with
% beamer compatibility, which has similar definitions but only up to 1. 
%    \begin{macrocode}
\expandafter\def\csname Parent2\endcsname{}
%    \end{macrocode}
% We need to disregard the columns macros introduced by the |beamer| class in the notes.  
%    \begin{macrocode}
\ifnotes%
  \renewenvironment{columns}{%
    \par\noindent%
    \begin{minipage}%
    \slidewidth\centering\leavevmode%
  }{%
    \end{minipage}\par\noindent%
  }%
  \newsavebox\columnbox%
  \renewenvironment{column}[1]{%
    \begin{lrbox}{\columnbox}\begin{minipage}{#1}%
  }{%
    \end{minipage}\end{lrbox}\usebox\columnbox%
  }%
\fi%
%    \end{macrocode}
%
%    \begin{macrocode}
\ifproblems%
  \newenvironment{problems}{}{}%
\else%
  \excludecomment{problems}%
\fi%
%</package>
%    \end{macrocode}
% \Finale
\endinput
% \iffalse
%%% Local Variables: 
%%% mode: doctex
%%% TeX-master: t
%%% End: 
% \fi
% LocalWords:  mikoslides dtx beamer omdoc notheorems noamsthm beamerarticle sc Licences
% LocalWords:  graphicx slidelabel stex amssymb tikz url CPERL  amsmath filedate
% LocalWords:  LoadClass RequirePackage DefRegister DefEnvironment omgroup rgb frameimage
% LocalWords:  afterDigestBegin setProperty LookupValue DefConstructor hyperref Tantau's
% LocalWords:  cmyk lec DefMacro titleslide ttitle RawTeX metadata etoolbox cls emph ugbc
% LocalWords:  noproblems linkcolor bookmarksopen citecolor urlcolor colorlinks Tantau di
% LocalWords:  breaklinks plainpages pdfpagelabels srcref iffalse texttt atwin exfig dt
% LocalWords:  mytwocolumns twintoo atwintoo scsys sc hbox raisebox cnxml impl notesfile
% LocalWords:  scshape ttfamily cnxlatex mathml activemath fileversion newpage tpp99 dd
% LocalWords:  maketitle tableofcontents ednote compactitem showmeta showmeta beamerclass
% LocalWords:  sectocframes sectocframes textwarning textwarning compactenum eq omgroups
% LocalWords:  includegraphics tracissue printbibliography textsf langle textsf rangle
% LocalWords:  langle ltxml metakeys newif ifnotes notesfalse ifsectocframes rangle putat
% LocalWords:  sectocframesfalse ifproblems problemstrue notestrue marginnote frontmatter
% LocalWords:  problemsfalse sectocframestrue mdframed noxcolor newcounter ifx equinymous
% LocalWords:  Hfootnote usetheme tikzinput usepgflibrary usetikzlibrary rhd ignorespaces
% LocalWords:  tikzmark textcomp newlength slidewidth setlength slidewidth miko setkeys
% LocalWords:  slideheight slideheight renewenvironment excludecomment itenize Gin@ewidth
% LocalWords:  slideframewidth slideframewidth surroundwithmdframed addmetakey Gen@ewidth
% LocalWords:  allowframebreaks allowdisplaybreaks metasetkeys stepcounter sbox mpart baz
% LocalWords:  currentlabel theslide renewcommand scriptsize scriptstyle hspace mchapter
% LocalWords:  medskip linewidth skipabove skipbelow frametitle newenvironment msection
% LocalWords:  slidelogoheight newsavebox slidelogo slidelogo jacobs-logo vbox ifcase
% LocalWords:  Attribuition-ShareAlike copyrightnotice footnotesize cclogo vss minipage
% LocalWords:  cclogo somerights ifcchref cchreffalse ifpackageloaded usebox mycgraphics
% LocalWords:  cchreftrue usebox newcommand hfill hfill definecolor definecolor endinput
% LocalWords:  defemph notemph stDMemph STpresent textcolor textwarigrening ltx HorIacJuc
% LocalWords:  pgfdeclareimage dbend pgfuseimage xspace titleframe titlepage mycgraphics
% LocalWords:  titleframewith hline vspace ttitlejoint newbox boxwith boxwith msubsection
% LocalWords:  putgraphicsat beginomgroup ifnum vfill vfill noindent leavevmode cscpnrr11
% LocalWords:  columnbox lrbox needwrapper unlist omtext bgroup autoclose pgf readXToken
% LocalWords:  mygraphics doctex NeedsTeXFormat textbackslash userdefinedwidth foobar
%  LocalWords:  includegrahics smalltextwarning msubsubsection msubsubsubsection ldots
%  LocalWords:  frameimages frameimages ifframeimages frameimagesfalse frameimagestrue
%  LocalWords:  expandafter csname endcsname specializes customization flexiformal colors
%  LocalWords:  initialize Initialization defindex realized itemize newrobustcmd
%  LocalWords:  specialize centering itemizations setsource setlicensing
